% =============================================================================
% Documento standalone: tabela de tempos de execucao das consultas.
% Compile com: pdflatex tabela_tempos.tex
% =============================================================================
\documentclass[12pt,a4paper]{article}
\usepackage[utf8]{inputenc}
\usepackage[T1]{fontenc}
\usepackage[brazil]{babel}
\usepackage[a4paper,top=2.5cm,bottom=2.5cm,left=2.5cm,right=2.5cm]{geometry}
\usepackage{float}
\usepackage{booktabs}
\usepackage{amsmath,amssymb}
\usepackage{caption}

\title{Avalia\c{c}\~ao de Desempenho --- Consultas SQL\\\large MATA60 -- Banco de Dados}
\author{Projeto Sistema de Eventos Acad\^emicos}
\date{}

\begin{document}
\maketitle

\section*{Metodologia}

Para cumprir o requisito \textbf{[Q2]} (avalia\c{c}\~ao de desempenho com pelo menos
10 execu\c{c}\~oes por consulta, m\'edia e desvio padr\~ao), cada uma das
30 consultas (10 intermedi\'arias e 20 avan\c{c}adas) foi executada
\textbf{10 vezes} no PostgreSQL~16 via \texttt{EXPLAIN (ANALYZE, TIMING ON)},
sendo coletado o \emph{Execution Time} reportado pelo planejador.
Antes de cada execu\c{c}\~ao foi emitido \texttt{DISCARD ALL} para limpar
caches de sess\~ao. Foram calculados, por consulta, a m\'edia ($\mu$),
o desvio padr\~ao amostral ($\sigma$), o m\'inimo e o m\'aximo dos tempos.

% Tabela gerada automaticamente por benchmark_results/aggregate.py
% Cada query foi executada 10 vezes via EXPLAIN ANALYZE no PostgreSQL 16.
\begin{table}[H]
\centering
\caption{Tempos de execu\c{c}\~ao das consultas (10 execu\c{c}\~oes cada via \texttt{EXPLAIN ANALYZE}).}
\label{tab:tempos-execucao}
\begin{tabular}{|c|r|r|r|r|r|}
\hline
\textbf{Consulta} & \textbf{N} & $\boldsymbol{\mu}$ \textbf{(ms)} & $\boldsymbol{\sigma}$ \textbf{(ms)} & \textbf{m\'in (ms)} & \textbf{m\'ax (ms)} \\
\hline
\multicolumn{6}{|c|}{\textbf{Consultas Intermedi\'arias}} \\
\hline
QI01 & 10 & 0.259 & 0.030 & 0.217 & 0.334 \\
QI02 & 10 & 1.293 & 0.037 & 1.221 & 1.355 \\
QI03 & 10 & 1.117 & 0.133 & 1.035 & 1.484 \\
QI04 & 10 & 2.662 & 0.130 & 2.524 & 2.930 \\
QI05 & 10 & 0.900 & 0.072 & 0.833 & 1.062 \\
QI06 & 10 & 0.732 & 0.021 & 0.684 & 0.753 \\
QI07 & 10 & 0.690 & 0.027 & 0.653 & 0.740 \\
QI08 & 10 & 0.614 & 0.043 & 0.534 & 0.668 \\
QI09 & 10 & 4.979 & 0.106 & 4.811 & 5.125 \\
QI10 & 10 & 2.779 & 0.041 & 2.706 & 2.834 \\
\hline
\multicolumn{6}{|c|}{\textbf{Consultas Avan\c{c}adas}} \\
\hline
QA01 & 10 & 3.054 & 0.084 & 2.941 & 3.177 \\
QA02 & 10 & 1.140 & 0.027 & 1.098 & 1.178 \\
QA03 & 10 & 2.870 & 0.072 & 2.792 & 2.986 \\
QA04 & 10 & 5.844 & 0.134 & 5.657 & 6.049 \\
QA05 & 10 & 2.079 & 0.055 & 2.004 & 2.169 \\
QA06 & 10 & 0.444 & 0.011 & 0.429 & 0.463 \\
QA07 & 10 & 5.272 & 0.062 & 5.200 & 5.400 \\
QA08 & 10 & 3.728 & 0.076 & 3.611 & 3.858 \\
QA09 & 10 & 2.817 & 0.091 & 2.694 & 2.985 \\
QA10 & 10 & 2.933 & 0.079 & 2.855 & 3.103 \\
QA11 & 10 & 3.778 & 0.039 & 3.721 & 3.852 \\
QA12 & 10 & 1.657 & 0.043 & 1.618 & 1.746 \\
QA13 & 10 & 1.350 & 0.020 & 1.325 & 1.387 \\
QA14 & 10 & 1.366 & 0.040 & 1.307 & 1.436 \\
QA15 & 10 & 1.326 & 0.017 & 1.305 & 1.352 \\
QA16 & 10 & 2.796 & 0.026 & 2.752 & 2.835 \\
QA17 & 10 & 6.381 & 0.093 & 6.259 & 6.572 \\
QA18 & 10 & 2.497 & 0.048 & 2.440 & 2.583 \\
QA19 & 10 & 0.667 & 0.017 & 0.639 & 0.692 \\
QA20 & 10 & 7.372 & 0.468 & 6.837 & 8.206 \\
\hline
\end{tabular}
\end{table}


\section*{Observa\c{c}\~oes}

\begin{itemize}
  \item Banco executado em container Docker (PostgreSQL 16) com o
        \emph{dataset} fornecido pelo grupo (\texttt{sql\_output/00\_executar\_tudo.sql}).
  \item Todas as 30 consultas executaram sem erro nas 10 repeti\c{c}\~oes ($N=10$).
  \item Os tempos consideram apenas o \emph{Execution Time} reportado por
        \texttt{EXPLAIN ANALYZE}, excluindo o tempo de planejamento.
  \item Os dados brutos est\~ao em \texttt{benchmark\_results/raw\_times.csv}
        e os agregados em \texttt{benchmark\_results/aggregated.csv}.
\end{itemize}

\end{document}
